\section{Twist families and the integer model}\label{sec:setup}

We first reduce upper bounds for arbitrary twists to a fixed finite
family of curves with coprime twisting parameters. This gives exact
conductor formulas while preserving the bounds for nonnegative sums
needed later. We then define the probability model for complete integer
residue averages.

We use unitary normalization:
\[
 L_F(s)=L(F,s+1/2)=\sum_{n\ge1}\lambda_F(n)n^{-s},
 \qquad q_F=\operatorname{cond}(F),
 \qquad C_F=\frac{\sqrt{q_F}}{2\pi}.
\]
The series converges absolutely in \(\Re s>1\).
Modularity and the local rules for the Hasse--Weil \(L\)-function give
real multiplicative coefficients and inverse Euler factors
\begin{equation}\label{eq:setup-euler}
 1-\lambda_F(p)p^{-s}+\xi_{F,p}p^{-2s},
 \qquad \xi_{F,p}=\mathbf 1_{p\nmid q_F}.
\end{equation}
At a good prime the two factor parameters have modulus one by Hasse's
bound \cite[Chapter~V, Theorem~2.3.1(a)]{Silverman2009}.
At a multiplicative prime there is one parameter, of modulus
\(p^{-1/2}\), and at an additive prime the factor is one
\cite[Definitions~17--18 and Theorem~19]{DokchitserAnni2014}.
In particular,
\begin{equation}\label{eq:setup-coefficients}
 |\lambda_F(n)|\le\tau(n).
\end{equation}
The completed function
\[
 \Lambda_F(s)=C_F^s\Gamma(s+1/2)L_F(s)
\]
is entire of finite order and satisfies
\begin{equation}\label{eq:setup-fe}
 \Lambda_F(s)=\epsilon_F\Lambda_F(1-s),\qquad \epsilon_F\in\{1,-1\}.
\end{equation}
Analytic continuation and the functional equation follow from
modularity \cite[Theorem~A]{BreuilConradDiamondTaylor2001}
and the standard completion
\cite[Appendix~C, Theorems~16.3--16.4]{Silverman2009}.
The standard finite-order bounds follow, for example, from the split
Mellin integral of the modular cusp form and its Fricke transform.
They also give polynomial growth in every fixed vertical strip for each
fixed curve \(F\), independently of the averaged estimates proved below.
Indeed, enclose the strip in \(a\le\Re s\le b\), with \(a\le-1\) and
\(b\ge2\). Absolute convergence bounds \(L_F\) on the right boundary;
\eqref{eq:setup-fe} and Stirling's formula give
\(L_F(a+it)\ll_{F,a}(1+|t|)^{1-2a}\) on the left boundary.
Take \(A=1-a\) and an integer \(N>1-2a\), so that
\(L_F(s)/(s+A)^N\) is bounded on both boundaries. For \(\varepsilon>0\) and \(0<\kappa<\pi/(b-a)\), multiply
this quotient by
\[
 \exp\left\{-\varepsilon\cos\left(\kappa
              \left(s-\frac{a+b}{2}\right)\right)\right\}.
\]
The multiplier has modulus at most one on the vertical boundaries and
at most \(\exp(-c\varepsilon\cosh(\kappa\Im s))\) on the horizontal
sides, for a fixed \(c>0\).
This dominates the finite-order growth. The maximum principle on
expanding rectangles, followed by \(\varepsilon\downarrow0\), therefore
gives \(L_F(\sigma+it)\ll_{F,a,b}(1+|t|)^N\) on the strip.
These individual bounds justify contour displacements. Their constants
may depend on the fixed curve; estimates uniform in the varying twist
parameters will be proved separately.

\subsection{A finite family closed under reduction}

Fix the primes
\[
 S=\{p:p\mid 2q_E\},\qquad Q_S=\prod_{p\in S}p.
\]
Let \(\mathcal F\) be the finite set, up to \(\Q\)-isomorphism, of curves
obtained from \(E\) by twisting by signed squarefree products of primes
in \(S\). The curve \(E\) is included. Every \(F\in\mathcal F\) has good
reduction outside \(S\).

Call a signed squarefree integer \(d\) \emph{admissible} if
\[
 d\equiv1\pmod4,\qquad (d,Q_S)=1.
\]
For a dyadic \(X\), put
\[
 \mathcal A(X)=\{d:d\text{ admissible},\ X\le|d|<2X\}.
\]

\begin{lemma}\label{setup:reduction}
For every \(F\in\mathcal F\) and signed squarefree \(b\), there are
\(F'\in\mathcal F\) and an admissible \(d\) such that
\[
 F^{(b)}\simeq (F')^{(d)},\qquad |d|\le |b|.
\]
The correspondence can be chosen to have bounded multiplicity, uniformly
in \(b\). Every sum of nonnegative functions of these twists can
therefore be bounded by a fixed multiple of the corresponding sums over
\(\mathcal F\) and admissible parameters of no greater height.
\end{lemma}

\begin{proof}
Write \(b=b_S b_0\), where \(b_S>0\) is the product of the primes of \(S\)
dividing \(b\). Then \(b_0\) is odd and prime to \(Q_S\).
Choose \(\varepsilon\in\{1,-1\}\) so that
\(d=\varepsilon b_0\equiv1\pmod4\), and put
\(F'=F^{(\varepsilon b_S)}\).
Quadratic twisting depends only on the square class and composes by
multiplication of square classes. Thus \(F'\) is in \(\mathcal F\) and
\((F')^{(d)}\simeq F^{(b)}\).
Moreover \(|d|=|b|/b_S\le |b|\).
There are at most \(2^{|S|+1}\) choices of the signed \(S\)-part.
Possible coincidences among the finitely many curves only change this
fixed multiplicity.
\end{proof}

\begin{lemma}\label{setup:local-factors}
For \(F\in\mathcal F\) and admissible \(d\),
\begin{equation}\label{eq:setup-twist}
 L_{F^{(d)}}(s)=\sum_{n\ge1}\lambda_F(n)\chi_d(n)n^{-s},
 \qquad q_{F^{(d)}}=q_F|d|^2,
 \qquad \chi_d(n)=\left(\frac dn\right).
\end{equation}
Here the symbol is the Kronecker symbol and \(d=1\) has the untwisted
convention. For an arbitrary signed squarefree \(b\), the same good-prime
twisting rule holds at primes outside \(S\): the local trace is multiplied
by \((b/p)\) when \(p\nmid b\), and the local factor is one when \(p\mid b\).
\end{lemma}

\begin{proof}
We compute the inverse local Euler factors using geometric Frobenius
on the inertia invariants of the dual of the rational Tate module,
and the conductor as the Artin conductor
\cite[Definitions~17--18 and Theorem~19]{DokchitserAnni2014}.
The rational Tate-module representation of a quadratic twist is the
original representation tensored with the corresponding quadratic
character. This follows from the isomorphism over the twisting extension,
whose conjugate differs by the scalar automorphism \([-1]\).

An unramified scalar twist preserves the conductor and multiplies
Frobenius eigenvalues on inertia invariants by the character value.
Since an admissible \(d\) is an odd fundamental discriminant, its character
is unramified at every prime of \(S\).
At a prime \(p\mid d\), the representation of \(F\) is unramified and the
quadratic character is tamely ramified. The tensor product has no inertia
invariants, Swan conductor zero, and Artin conductor exponent two.
At all other primes the scalar twist is unramified.
Computing the inertia-invariant Euler factors and the conductor gives
\eqref{eq:setup-twist}.
For an arbitrary squarefree \(b\), these same computations apply at odd
primes outside \(S\), regardless of the character's behavior at two.
\end{proof}

The constants in what follows are chosen uniformly over the finite set
\(\mathcal F\).
The nonincreasing height in Lemma~\ref{setup:reduction} is useful in the
second-moment induction: a dual squarefree parameter of absolute value
at most \(M/2\) remains in a smaller dyadic block after reduction.

\subsection{Complete integer averages}

Fix \(c\in\{1,-1\}\). For \(m\in\Z\), define a completely multiplicative
formal character on positive integers by
\[
 z_m(2)=c,\qquad z_m(p)=\left(\frac mp\right)
 \quad(p\text{ odd}).
\]
Let \(Z\) be the completely multiplicative random model with \(Z(2)=c\)
and independent odd-prime values distributed by
\begin{equation}\label{eq:setup-model}
 \Prob(Z(p)=0)=\frac1p,\qquad
 \Prob(Z(p)=1)=\Prob(Z(p)=-1)=\frac{1-1/p}{2}.
\end{equation}
For every polynomial involving only finitely many prime values,
its expectation is the average over a complete residue system for the
corresponding odd-prime periods. This follows from the Chinese remainder
theorem and the equal numbers of nonzero quadratic residues and
nonresidues modulo each odd prime.

We use this model for averages over all integers \(m\).
Nonnegative estimates can subsequently be restricted to admissible
squarefree \(d\), separated into the two classes \(\chi_d(2)=c\);
on either class, \(z_d(n)=\chi_d(n)\) for every \(n\).
